\section{Architecture}


\subsection{Server Factory}

\begin{figure}[t]
\centering
\begin{tikzpicture}[
    node distance=0.8cm,
    box/.style={rectangle, draw, minimum width=2.5cm, minimum height=0.7cm, align=center, font=\small},
    arrow/.style={->, thick}
]
\node[box, fill=blue!10] (server) {MCP Server};
\node[box, below=of server, fill=green!10] (registry) {Tool Registry};
\node[box, below left=0.5cm and -0.5cm of registry] (file) {File Ops};
\node[box, below=of registry] (shell) {Shell};
\node[box, below right=0.5cm and -0.5cm of registry] (search) {Search};
\node[box, left=of file] (edit) {Edit};
\node[box, right=of search] (ui) {UI};
\node[box, below=of shell] (orch) {Orchestration};

\draw[arrow] (server) -- (registry);
\draw[arrow] (registry) -- (file);
\draw[arrow] (registry) -- (shell);
\draw[arrow] (registry) -- (search);
\draw[arrow] (registry) -- (edit);
\draw[arrow] (registry) -- (ui);
\draw[arrow] (registry) -- (orch);
\end{tikzpicture}
\caption{MCP Server architecture with tool registry and categories.}
\label{fig:arch}
\end{figure}

The server is created via a factory function (Figure~\ref{fig:arch}):

\begin{lstlisting}[language=JavaScript]
export async function createMCPServer(config?: {
  name?: string;
  version?: string;
  projectPath?: string;
  customTools?: Tool[];
  toolConfig?: ToolConfig;
}) {
  const configuredTools = getConfiguredTools({
    ...toolConfig,
    customTools: [...(toolConfig.customTools || []),
                   ...customTools]
  });

  const server = new Server(
    { name, version },
    { capabilities: { tools: {}, resources: {} } }
  );

  server.setRequestHandler(ListToolsRequestSchema,
    async () => ({
      tools: configuredTools.map(tool => ({
        name: tool.name,
        description: tool.description,
        inputSchema: tool.inputSchema
      }))
    }));

  return { server, tools: configuredTools, start,
           addTool, removeTool };
}
\end{lstlisting}

\subsection{Tool Configuration}

The \texttt{ToolConfig} interface enables selective tool loading:

\begin{lstlisting}[language=TypeScript]
export interface ToolConfig {
  enableCore?: boolean;      // File, shell, edit
  enableUI?: boolean;        // Playwright, UI registry
  enableAutoGUI?: boolean;   // PyAutoGUI adapters
  enableOrchestration?: boolean;  // Agent tools
  enableUIRegistry?: boolean;
  enableGitHubUI?: boolean;
  enabledCategories?: string[];
  disabledTools?: string[];
  customTools?: Tool[];
}
\end{lstlisting}

This design allows deployments to enable only required tools, reducing attack surface and memory footprint.

\subsection{Tool Categories}

\subsubsection{File Operations}

Four tools for filesystem interaction:
\begin{itemize}
    \item \texttt{read\_file}: Read file contents with optional line range
    \item \texttt{write\_file}: Write content to file (create or overwrite)
    \item \texttt{list\_directory}: List directory contents with metadata
    \item \texttt{file\_info}: Get file metadata (size, permissions, timestamps)
\end{itemize}

All file operations validate paths against a configurable allowlist to prevent directory traversal attacks.

\subsubsection{Shell Execution}

\begin{lstlisting}[language=TypeScript]
interface ShellToolInput {
  command: string;
  cwd?: string;
  timeout?: number;      // Default: 30000ms
  background?: boolean;  // Run asynchronously
}
\end{lstlisting}

The shell tool executes commands in isolated processes with:
\begin{itemize}
    \item Configurable timeouts to prevent runaway processes
    \item Background execution for long-running commands
    \item Output streaming for real-time feedback
    \item Exit code propagation for error handling
\end{itemize}

\subsubsection{Edit Operations}

Precise file editing with three modes:
\begin{itemize}
    \item \texttt{edit\_replace}: Replace exact string matches
    \item \texttt{edit\_insert}: Insert at line number
    \item \texttt{edit\_delete}: Delete line range
\end{itemize}

The edit tool uses fuzzy matching to handle whitespace variations while preserving indentation.

\subsubsection{Search Engine}

\begin{figure}[t]
\centering
\begin{tikzpicture}[
    node distance=0.6cm,
    box/.style={rectangle, draw, minimum width=2cm, minimum height=0.6cm, align=center, font=\small},
    arrow/.style={->, thick}
]
\node[box, fill=yellow!20] (query) {Query};
\node[box, below left=0.8cm and 0cm of query] (text) {Text Search};
\node[box, below=0.8cm of query] (ast) {AST Search};
\node[box, below right=0.8cm and 0cm of query] (vector) {Vector Search};
\node[box, below=1.5cm of ast, fill=green!20] (merge) {Result Merger};
\node[box, below=of merge, fill=blue!20] (rank) {Ranked Results};

\draw[arrow] (query) -- (text);
\draw[arrow] (query) -- (ast);
\draw[arrow] (query) -- (vector);
\draw[arrow] (text) -- (merge);
\draw[arrow] (ast) -- (merge);
\draw[arrow] (vector) -- (merge);
\draw[arrow] (merge) -- (rank);
\end{tikzpicture}
\caption{Unified search engine with three retrieval strategies.}
\label{fig:search}
\end{figure}

The unified search engine (Figure~\ref{fig:search}) combines:

\paragraph{Text Search.} Powered by ripgrep for high-performance regex matching across large codebases. Supports glob filtering, context lines, and match highlighting.

\paragraph{AST Search.} Uses tree-sitter grammars to parse source files and match structural patterns. Enables queries like ``find all async functions returning Promise$<$T$>$''.

\paragraph{Vector Search.} Embeds code snippets using a local model and stores vectors in LanceDB for similarity search. Useful for semantic queries like ``find authentication logic''.

Results are merged using reciprocal rank fusion:
\begin{equation}
    \text{score}(d) = \sum_{s \in \text{strategies}} \frac{1}{k + \text{rank}_s(d)}
\end{equation}
where $k$ is a smoothing constant (default: 60).

\subsubsection{UI Automation}

Playwright integration enables cross-browser automation:

\begin{lstlisting}[language=TypeScript]
interface UIToolInput {
  action: 'navigate' | 'click' | 'type' |
          'screenshot' | 'evaluate';
  selector?: string;
  url?: string;
  text?: string;
  code?: string;  // For evaluate
  device?: 'mobile' | 'tablet' | 'desktop';
}
\end{lstlisting}

Features include:
\begin{itemize}
    \item Device emulation (iPhone, Pixel, iPad, desktop)
    \item Network interception for mocking APIs
    \item Trace recording for debugging
    \item PDF generation for documentation
\end{itemize}

\subsubsection{Agent Orchestration}

Tools for multi-agent coordination:
\begin{itemize}
    \item \texttt{spawn\_agent}: Create sub-agent with specific role
    \item \texttt{send\_message}: Inter-agent communication
    \item \texttt{await\_result}: Wait for agent completion
    \item \texttt{terminate\_agent}: Clean shutdown
\end{itemize}

